% ======================================================
% 第 71 天
% ======================================================
\setday{71}{重积分的应用}

\oneproblem{
    求曲面 $x^2+y^2=az$ 和 $z=2a-\sqrt{x^2+y^2}\ (a>0)$ 所围立体的表面积.
}

\oneproblem{
    设 $a,b$ 为实数，函数 $z=2+ax^2+by^2$ 在点 $(3,4)$ 处的方向导数中，沿方向 $\vec{l}=-3\vec{i}-4\vec{j}$ 的方向导数最大，最大值为 $10$.
    \begin{enumerate}[label=(\arabic*)]
        \item 求 $a,b$；
        \item 求曲面 $z=2+ax^2+by^2\ (z\ge0)$ 的面积.
    \end{enumerate}
}

\oneproblem{
    设半径为 $R$ 的球面 $\Sigma$ 的球心在定球面 $x^2+y^2+z^2=a^2\ (a>0)$ 上，问当 $R$ 取何值时，球面 $\Sigma$ 在定球面内部的那部分的面积最大？
}

\oneproblem{
    设有一高度为 $h(t)$ ($t$ 为时间) 的雪堆在融化过程，其侧面满足方程 $z=h(t)-\dfrac{2(x^2+y^2)}{h(t)}$。设长度单位为厘米，时间单位为小时。已知体积减少的速率与侧面积成正比（比例系数为 $0.9$），问高度为 $130$（厘米）的雪堆全部融化需多少小时？
}

\oneproblem{
    设曲线 $L$ 的方程为 $y=\dfrac{1}{4}x^2-\dfrac{1}{2}\ln x\ (1\le x\le e)$.
    \begin{enumerate}[label=(\arabic*)]
        \item 求 $L$ 的弧长；
        \item 设 $D$ 是由曲线 $L$，直线 $x=1,x=e$ 及 $x$ 轴所围成的平面图形，求 $D$ 的形心的横坐标.
    \end{enumerate}
}

\oneproblem{
    设有一半径为 $R$ 的球体，$P_0$ 是此球的表面上的一个定点，球体上任一点的密度与该点到 $P_0$ 距离的平方成正比（比例常数 $k>0$），求球体的重心位置.
}

\oneproblem{
    设直线 $L$ 过 $A(1,0,0),B(0,1,1)$ 两点，将 $L$ 绕 $z$ 轴旋转一周得到曲面 $\Sigma$，$\Sigma$ 与平面 $z=0,z=2$ 所围成的立体为 $\Omega$.
    \begin{enumerate}[label=(\arabic*)]
        \item 求曲面 $\Sigma$ 的方程；
        \item 求 $\Omega$ 的形心坐标.
    \end{enumerate}
}

\oneproblem{
    设 $\Omega$ 是由锥面 $x^2+(y-z)^2=(1-z)^2\ (0\le z\le1)$ 与平面 $z=0$ 所围成的锥体，求 $\Omega$ 的形心坐标.
}

\oneproblem{
    设 $l$ 是过原点、方向为 $(\alpha,\beta,\gamma)$（其中 $\alpha^2+\beta^2+\gamma^2=1$）的直线，均匀椭球 $\dfrac{x^2}{a^2}+\dfrac{y^2}{b^2}+\dfrac{z^2}{c^2}\le1$（其中 $0<c<b<a$，密度为 $1$）绕 $l$ 旋转.
    \begin{enumerate}[label=(\arabic*)]
        \item 求其转动惯量；
        \item 求其转动惯量关于方向 $(\alpha,\beta,\gamma)$ 的最大值和最小值.
    \end{enumerate}
}

\oneproblem{
    求高为 $H$、底半径为 $R$ 且密度均匀的圆柱体，对圆柱底面中心一单位质点的引力.
}

\oneproblem{
    设 $D$ 为椭圆 $\dfrac{x^2}{a^2}+\dfrac{y^2}{b^2}\le a\ (a>b>0)$，面密度为 $\rho$ 的均质薄板；$l$ 为通过椭圆焦点 $(-c,0)$（其中 $c^2=a^2-b^2$）垂直于薄板的旋转轴.
    \begin{enumerate}[label=(\arabic*)]
        \item 求薄板 $D$ 绕 $l$ 旋转的转动惯量 $J$；
        \item 对于固定的转动惯量，讨论椭圆薄板的面积是否有最大值和最小值.
    \end{enumerate}
}

\oneproblem{
    在研究某流行病的传播时，假设患病者将疾病传染给健康者的概率是两者之间距离的函数。现考虑一个半径为 $10$ km 的圆形城市（所占平面区域用 $D$ 表示）。假设城中的人口和患病者都是均匀分布的（每 km$^2$ 有 $k$ 个患者），又设城中某点 $M$ 处的健康者被位于点 $P$ 处的患者传染的概率为 $f(P)=1-\dfrac{1}{20}d(P,M)$，其中 $d(P,M)$ 表示点 $P$ 到点 $M$ 的距离。假设位于点 $M$ 处的健康者传染上该病的概率（称为被感染率，记作 $E(M)$）等于城市中所有患者将疾病传染给他的概率之和.
    \begin{enumerate}[label=(\arabic*)]
        \item 试用二重积分表示 $E(M)$；
        \item 设 $A$ 点位于城市中心，$B$ 点位于城市边缘，试计算 $E(A)$ 和 $E(B)$ 的值；并问居住在城市中心的居民与居住在城市边缘处的居民相比，谁的被感染率大？
    \end{enumerate}
}

% ======================================================
% 第 72 天
% ======================================================
\setday{72}{含参变量的积分}

\oneproblem{
    求极限 $\displaystyle \lim_{\alpha\to0} \int_0^{1+\alpha} \dfrac{dx}{1+x^2+\alpha^2}$.
}

\oneproblem{
    设函数 $f(x)$ 是 $[0,1]$ 上连续的正值函数，试研究函数 $I(y)=\displaystyle \int_0^1 \dfrac{yf(x)}{x^2+y^2}\,dx$ 的连续性.
}

\oneproblem{
    求 $I(y)=\displaystyle \int_y^{y^2} \dfrac{\cos xy}{x}\,dx$ 的导数 $I'(y)$.
}

\oneproblem{
    设 $F(x,y)=\displaystyle \int_{\frac{y}{x}}^{xy} (x-yz)f(z)\,dz$，其中 $f$ 可微，求 $F_{xy}(x,y)$.
}

\oneproblem{
    求积分 $\displaystyle \int_0^1 \sin\left(\ln\dfrac{1}{x}\right)\dfrac{x^b-x^a}{\ln x}\,dx,\ b>a>0$.
}

\oneproblem{
    利用含参积分求导方法求积分 $I(\alpha)=\displaystyle \int_0^\pi \ln(1-2\alpha\cos x+\alpha^2)\,dx$.
}

\oneproblem{
    利用积分号下求导法求积分 $I(a)=\displaystyle \int_0^{\frac{\pi}{2}} \dfrac{\arctan(a\tan x)}{\tan x}\,dx,\ |a|<1$.
}

\oneproblem{
    设 $f(x)$ 在 $[a,b]$ 连续，证明：$y(x)=\dfrac{1}{k}\int_c^x f(t)\sin k(x-t)\,dt$（其中 $a,c\in[a,b]$）满足微分方程 $y''+k^2y=f(x)$.
}

\oneproblem{
    设 $f(x)$ 为连续函数，$F(x)=\displaystyle \int_0^h \left[\int_0^h f(x+\xi+\eta)\,d\eta\right]d\xi$，求 $F''(x)$.
}

\oneproblem{
    讨论下列反常积分的一致收敛性：
    \begin{enumerate}[label=(\arabic*)]
        \item $\displaystyle \int_1^{+\infty} \dfrac{y^2-x^2}{(x^2+y^2)^2}\,dx,\ y\in(-\infty,+\infty)$.
        \item $\displaystyle \int_0^{+\infty} e^{-xy}\dfrac{\sin x}{x}\,dx,\ y\in[0,+\infty)$.
    \end{enumerate}
}

\oneproblem{
    证明：
    \begin{enumerate}[label=(\arabic*)]
        \item $\displaystyle \int_1^{+\infty} e^{-yx^2}\sin y\,dx$ 关于 $y\in[0,+\infty)$ 一致收敛；
        \item $\displaystyle \int_1^{+\infty} e^{-yx^2}\sin y\,dy$ 关于 $x\in[0,+\infty)$ 不一致收敛.
    \end{enumerate}
}

\oneproblem{
    证明函数 $F(\alpha)=\displaystyle \int_1^{+\infty} \dfrac{\cos x}{x^\alpha}\,dx$ 在 $\alpha>0$ 内连续.
}

\oneproblem{
    计算 $g(\alpha)=\displaystyle \int_1^{+\infty} \dfrac{\arctan\alpha x}{x^2\sqrt{x^2-1}}\,dx$.
}

\oneproblem{
    利用积分号下求导求积分 $I_n(a)=\displaystyle \int_0^{+\infty} \dfrac{dx}{(x^2+a)^{n+1}}$，$n$ 为正整数，$a>0$.
}

\oneproblem{
    计算积分 $\displaystyle \int_0^{+\infty} e^{-\left(x^2+\frac{a^2}{x^2}\right)}\,dx$.
}

\oneproblem{
    \begin{enumerate}[label=(\arabic*)]
        \item 设函数 $f(x)$ 在 $[0,+\infty)$ 上连续，且 $\displaystyle \lim_{x\to+\infty} f(x)=0$。证明 $\displaystyle \int_0^{+\infty} \dfrac{f(ax)-f(bx)}{x}\,dx=f(0)\ln\dfrac{b}{a},\ (a,b>0)$.
        \item 设 $\beta>\alpha>0$，计算 $\displaystyle \int_0^{+\infty} \dfrac{e^{-\alpha x^2}-e^{-\beta x^2}}{x^2}\,dx$.
    \end{enumerate}
}

\oneproblem{
    计算下列积分：
    \begin{enumerate}[label=(\arabic*),itemsep = 4cm]
        \item $\displaystyle \int_0^1 \dfrac{dx}{\sqrt[n]{1-x^n}}$.
        \item $\displaystyle \int_0^{\frac{\pi}{2}} \sin^7x\cos^{\frac{1}{2}}x\,dx$.
        \item $\displaystyle \int_0^1 \ln\Gamma(x)\,dx$.
    \end{enumerate}
}

% ======================================================
% 第 73 天
% ======================================================
\setday{73}{对弧长的曲线积分}

\oneproblem{
    设平面曲线 $L$ 为下半圆 $y=-\sqrt{1-x^2}$，计算对弧长的曲线积分 $\displaystyle \int_L (x^2+y^2)\,ds$.
}

\oneproblem{
    计算 $\displaystyle \int_L |y|\,ds,\ L$ 为双纽线 $(x^2+y^2)^2=a^2(x^2-y^2)\ (a>0)$.
}

\oneproblem{
    计算 $\displaystyle \int_L \sqrt{2y^2+z^2}\,ds$，其中 $L$ 为 $x^2+y^2+z^2=a^2$ 与 $y=x$ 的交线.
}

\oneproblem{
    设 $L$ 为椭圆 $\dfrac{x^2}{4}+\dfrac{y^2}{3}=1$，其周长记为 $a$，计算对弧长的曲线积分 $\displaystyle \oint_L (2xy+3x^2+4y^2)\,ds$.
}

\oneproblem{
    记空间曲线 $\Gamma:\begin{cases} x^2+y^2+z^2=a^2 \\ x+y+z=0 \end{cases}\ (a>0)$，计算对弧长的曲线积分 $\displaystyle \oint_\Gamma (1+x)^2\,ds$.
}

\oneproblem{
    求八分之一球面 $x^2+y^2+z^2=R^2,\ x\ge0,y\ge0,z\ge0$ 的边界曲线的重心，设曲线的线密度 $\rho=1$.
}

\oneproblem{
    曲线 $S$ 由 $x^2+y^2+z^2=1$ 与 $x+y+z=0$ 相交而成，计算对弧长的曲线积分 $\displaystyle \oint_S xy\,ds$.
}

\oneproblem{
    设曲线 $C:x^2+xy+y^2=a^2$ 的长度为 $L$，求 $I=\displaystyle \int_C \dfrac{a\sin(e^x)+b\sin(e^y)}{\sin(e^x)+\sin(e^y)}\,ds$.
}

\oneproblem{
    计算 $I=\displaystyle \oint_\Gamma (x^2+z^2)\,ds$，其中 $\Gamma$ 为球面 $(x-1)^2+(y+1)^2+z^2=a^2\ (a>0)$ 和平面 $x+y+z=0$ 的交线.
}

\oneproblem{
    设圆柱螺线 $x=\cos t,y=\sin t,z=t\ (0\le t\le\dfrac{\pi}{2})$ 的密度分布与 $x,y$ 无关而与 $z$ 成正比，求这一段螺线的质量和质心.
}

\oneproblem{
    设函数 $f(x)$ 为连续函数，$\Omega$ 为曲面 $\Sigma_1:z=x^2+y^2$ 与 $\Sigma_2:z=t\ (t>0)$ 所围成的立体，$L$ 为曲面 $\Sigma_1$ 与 $\Sigma_2$ 的交线，已知对任意实数 $t>0$，都有 $\displaystyle \iiint_\Omega f(z)\,dV=\pi f(t)+\displaystyle \oint_L (x^2+y^2)^{\frac{3}{2}}\,ds$，求函数 $f(x)$ 的表达式.
}

\oneproblem{
    试证明曲线 $(x-y)^2=3(x+y),\ x^2-y^2=\dfrac{9}{8}z^2$ 上从 $O(0,0,0)$ 到 $P(x_0,y_0,z_0)$ 的一段曲线 $\Gamma$ 的弧长仅仅与 $z_0$ 的取值有关而与 $x_0,y_0$ 无关，并计算当 $z=1$ 时的弧长.
}

% ======================================================
% 第 74 天
% ======================================================
\setday{74}{对坐标的曲线积分}

\oneproblem{
    计算曲线积分 $\displaystyle \int_L \sin2x\,dx+2(x^2-1)y\,dy$，其中 $L$ 是曲线 $y=\sin x$ 上从点 $(0,0)$ 到点 $(\pi,0)$ 的一段.
}

\oneproblem{
    设 $\Gamma$ 为空间曲线 $\begin{cases} x=\pi\sin(t/2) \\ y=t-\sin t \\ z=\sin2t \end{cases}$ 从 $t=0$ 到 $t=\pi$ 的一段，计算第二型曲线积分 $\displaystyle \int_\Gamma e^{\sin x}(\cos x\cos y\,dx-\sin y\,dy)+\cos z\,dz$.
}

\oneproblem{
    设 $L$ 是柱面方程 $x^2+y^2=1$ 与平面 $z=x+y$ 的交线，从 $z$ 轴正向往 $z$ 轴负向看去为逆时针方向，计算对坐标的曲线积分 $\displaystyle \oint_L xz\,dx+x\,dy+\dfrac{y^2}{2}\,dz$.
}

\oneproblem{
    质点 $P$ 沿着以 $A,B$ 为直径的半圆周，从点 $A(1,2)$ 运动到点 $B(3,4)$ 的过程中受变力 $\vec{F}$ 作用（见图），$\vec{F}$ 的大小等于点 $P$ 与原点 $O$ 之间的距离，其方向垂直于线段 $OP$ 且与 $y$ 轴正向的夹角小于 $\dfrac{\pi}{2}$，求变力 $\vec{F}$ 对质点 $P$ 所作的功.
}

\oneproblem{
    在过点 $O(0,0)$ 和 $A(\pi,0)$ 的曲线族 $y=a\sin x\ (a>0)$ 中，求一条曲线 $L$，使沿该曲线从 $O$ 到 $A$ 的积分 $\displaystyle \int_L (1+y^3)\,dx+(2x+y)\,dy$ 的值最小.
}

\oneproblem{
    在变力 $\vec{F}=yz\vec{i}+zx\vec{j}+xy\vec{k}$ 的作用下，质点由原点沿直线运动到椭球面 $\dfrac{x^2}{a^2}+\dfrac{y^2}{b^2}+\dfrac{z^2}{c^2}=1$ 上第一卦限的点 $M(\xi,\eta,\zeta)$，问当 $\xi,\eta,\zeta$ 取何值时，力 $\vec{F}$ 所作的功 $W$ 最大？并求出 $W$ 的最大值.
}

\oneproblem{
    已知平面区域 $D=\{(x,y)|0\le x\le\pi,0\le y\le\pi\}$，$L$ 为 $D$ 的正向边界。试证：
    \begin{enumerate}[label=(\arabic*),itemsep = 4cm]
        \item $\displaystyle \oint_L xe^{\sin y}\,dy-ye^{-\sin x}\,dx=\oint_L xe^{-\sin y}\,dy-ye^{\sin x}\,dx$.
        \item $\displaystyle \oint_L xe^{\sin y}\,dy-ye^{-\sin x}\,dx\ge2\pi^2$.
        \item $\displaystyle \oint_L xe^{\sin y}\,dy-ye^{-\sin x}\,dx\ge\dfrac{5}{2}\pi^2$.
    \end{enumerate}
}

\oneproblem{
    设 $M=\max_{(x,y)\in C}\sqrt{P^2+Q^2}$，$P(x,y),Q(x,y)$ 在曲线 $C$ 上连续，曲线段 $C$ 的长度为 $L$\\证明：$\left|\displaystyle \int_C P\,dx+Q\,dy\right|\le M\cdot L$。并利用该不等式估计积分 $I_R=\displaystyle \oint_{C_R} \dfrac{(y-1)\,dx+(x+1)\,dy}{(x^2+y^2+2x-2y+2)^2}$，其中 $C_R$ 为圆周 $(x+1)^2+(y-1)^2=R^2$ 的正向，并求 $\displaystyle \lim_{R\to+\infty} |I_R|$.
}

\oneproblem{
    设 $\Gamma$ 为曲线 $x=t,y=t^2,z=t^3$ 上相应于 $t$ 从 $0$ 变到 $1$ 的曲线弧，把对坐标的曲线积分 $I=\displaystyle \int_\Gamma dx+2x\,dy+3y\,dz$ 化成对弧长的曲线积分，并分别利用两种的计算方法计算积分值.
}

\oneproblem{
    计算 $I=\displaystyle \oint_\Gamma |\sqrt{3}y-x|\,dx-5z\,dz$，其中 $\Gamma:\begin{cases} x^2+y^2+z^2=8, \\ x^2+y^2=2z \end{cases}$ 从 $z$ 轴正向往坐标原点看去取逆时针方向.
}

\oneproblem{
    设 $I_a(r)=\displaystyle \int_C \dfrac{y\,dx-x\,dy}{(x^2+y^2)^a}$，其中 $a$ 为常数，曲线 $C$ 为椭圆 $x^2+xy+y^2=r^2$，取正向，求极限 $\displaystyle \lim_{r\to+\infty} I_a(r)$.
}

% ======================================================
% 第 75 天
% ======================================================
\setday{75}{格林公式及其应用}

\oneproblem{
    设 $L_1:x^2+y^2=1,\ L_2:x^2+y^2=2,\ L_3:x^2+2y^2=2,\ L_4:2x^2+y^2=2$ 为四条逆时针方向的平面曲线，记 $I_i=\displaystyle \oint_{L_i} \left(y+\dfrac{y^3}{6}\right)dx+\left(2x-\dfrac{x^3}{3}\right)dy\ (i=1,2,3,4)$，则 $\max\{I_1,I_2,I_3,I_4\}=$（ \quad ）
    \begin{multicols}{2}
    \begin{enumerate}[label=(\Alph*)]
        \item $I_1$
        \item $I_2$
        \item $I_3$
        \item $I_4$
    \end{enumerate}
    \end{multicols}
}

\oneproblem{
    设曲线积分 $I=\displaystyle \oint_L \dfrac{xdy-ydx}{|x|+|y|}$，其中 $L$ 是以 $(1,0),(0,1),(-1,0),(0,-1)$ 为顶点的正方形的边界曲线，方向为逆时针方向，则 $I=$ \underline{\hbox to 20mm{}}.
}

\oneproblem{
    设 $L$ 为沿圆 $x^2+y^2=a^2\ (a>0)$ 从 $(0,a)$ 依逆时针到点 $(0,-a)$ 的半圆，计算积分
    $
    I=\int_L \dfrac{y^2}{\sqrt{a^2+x^2}}dx+\left[ax+2y\ln(x+\sqrt{a^2+x^2})\right]dy.
    $
}

\oneproblem{
    设 $L$ 为 $x^2+y^2=2x\ (y\ge0)$ 上从 $O(0,0)$ 到 $A(2,0)$ 的一段弧，连续函数 $f(x)$ 满足
    $
    f(x)=x^2+\int_L y\left[f(x)+e^x\right]dx+(e^x-xy^2)dy,
    $
    求 $f(x)$.
}

\oneproblem{
    设在上半平面 $D=\{(x,y)|y>0\}$ 内，数 $f(x,y)$ 是有连续偏导数，且对任意的 $t>0$ 都有 $f(tx,ty)=t^{-2}f(x,y)$。证明：对 $D$ 内的任意分段光滑的有向简单闭曲线 $L$，都有
    $
    \oint_L yf(x,y)dx-xf(x,y)dy=0.
    $
}

\oneproblem{
    证明：若 $L$ 为平面上封闭曲线，$\vec{m}$ 为任意方向向量，则 $\displaystyle \oint_L \cos(\vec{m},\vec{n})ds=0$，其中 $\vec{n}$ 为曲线 $L$ 的外法线方向.
}

\oneproblem{
    设 $C$ 是圆周 $(x-1)^2+(y-1)^2=1$ 的边界曲线的正向，$f(x)$ 为大于零的连续函数。证明：
    $
    \oint_C xf(y)dy-\dfrac{y}{f(x)}dx\ge2\pi.
    $
}

\oneproblem{
    设 $D$ 是由极坐标方程曲线 $C:r=1+\cos\theta$ 所围成的闭区域，面积为 $A$。曲线 $C$ 的方向为逆时针方向，函数 $u=u(x,y)$ 在 $D$ 上具有二阶连续偏导数，且 $u''_{xx}+u''_{yy}=1$。证明：
    $
    \oint_C \dfrac{\partial u}{\partial n}ds=A,
    $
    其中 $\dfrac{\partial u}{\partial n}$ 是 $u$ 沿 $D$ 的边界外向法线的方向导数，并求此积分值.
}

\oneproblem{
    设函数 $f(x,y)$ 在区域 $D:x^2+y^2\le1$ 上有二阶连续偏导数，且 $\dfrac{\partial^2 f}{\partial x^2}+\dfrac{\partial^2 f}{\partial y^2}=e^{-(x^2+y^2)}$。证明：
    $
    \iint_D \left(x\dfrac{\partial f}{\partial x}+y\dfrac{\partial f}{\partial y}\right)\dd x \dd y=\dfrac{\pi}{2e}.
    $
}

\oneproblem{
    设函数 $f(x,y)$ 在区域 $D=\{(x,y)|x^2+y^2\le a^2\}$ 上具有一阶连续偏导数，且满足 $f(x,y)|_{x^2+y^2=a^2}=a^2$ 以及 $\max_{(x,y)\in D}\left[\left(\dfrac{\partial f}{\partial x}\right)^2+\left(\dfrac{\partial f}{\partial y}\right)^2\right]=a^2$，其中 $a>0$。证明：
    $
    \left|\iint_D f(x,y)\dd x \dd y\right|\le\dfrac{4}{3}\pi a^4.
    $
}

\oneproblem{
    设连续可微函数 $z=z(x,y)$ 由方程 $F(xz-y,x-yz)=0$（其中 $F(u,v)$ 有连续偏导数）唯一确定，$L$ 为正向单位圆周，试求：
    $
    I=\oint_L (xz^2+2yz)dy-(2xz+yz^2)dx.
    $
}

\oneproblem{
    设 $f(x,y)$ 是 $D=\{(x,y)|x^2+y^2\le1\}$ 上二次连续可微函数，满足 $\dfrac{\partial^2 f}{\partial x^2}+\dfrac{\partial^2 f}{\partial y^2}=x^2y^2$，计算积分
    \[
    I=\iint_D \left(\dfrac{x}{\sqrt{x^2+y^2}}\dfrac{\partial f}{\partial x}+\dfrac{y}{\sqrt{x^2+y^2}}\dfrac{\partial f}{\partial y}\right)\dd x \dd y.
    \]
}

\oneproblem{
    设函数 $f(x,y)$ 在闭区域 $D=\{(x,y)|x^2+y^2\le1\}$ 上具有二阶连续偏导数，且 $\dfrac{\partial^2 f}{\partial x^2}+\dfrac{\partial^2 f}{\partial y^2}=x^2+y^2$，求
    \[
    \lim_{r\to0^+}\dfrac{\displaystyle \iint_{x^2+y^2\le r^2} \left(x\dfrac{\partial f}{\partial x}+y\dfrac{\partial f}{\partial y}\right)\dd x \dd y}{(\tan r-\sin r)^2}.
    \]
}

% ======================================================
% 第 76 天
% ======================================================
\setday{76}{积分与路径无关的等价描述及其应用}

\oneproblem{
    设曲线 $L:f(x,y)=1$（$f(x,y)$ 具有一阶连续偏导数），过第 2 象限内的点 $M$ 和第 4 象限内的点 $N$，$\Gamma$ 为 $L$ 上从点 $M$ 到点 $N$ 的一段弧，则下列积分小于零的是（ \quad ）
    \begin{multicols}{2}
    \begin{enumerate}[label=(\Alph*)]
        \item $\displaystyle \int_\Gamma f(x,y)dx$
        \item $\displaystyle \int_\Gamma f(x,y)dy$
        \item $\displaystyle \int_\Gamma f(x,y)ds$
        \item $\displaystyle \int_\Gamma f'_x(x,y)dx+f'_y(x,y)dy$
    \end{enumerate}
    \end{multicols}
}

\oneproblem{
    设函数 $Q(x,y)=\dfrac{x}{y^2}$。如果对上半平面 $(y>0)$ 内的任意有向光滑闭曲线 $C$ 都有 $\displaystyle \oint_C P(x,y)dx+Q(x,y)dy=0$，那么函数 $P(x,y)$ 可取为（ \quad ）
    \begin{multicols}{2}
    \begin{enumerate}[label=(\Alph*)]
        \item $y-\dfrac{x^2}{y^3}$
        \item $\dfrac{1}{y}-\dfrac{x^2}{y^3}$
        \item $\dfrac{1}{x}-\dfrac{1}{y}$
        \item $x-\dfrac{1}{y}$
    \end{enumerate}
    \end{multicols}
}

\oneproblem{
    已知 $du(x,y)=\dfrac{ydx-xdy}{3x^2-2xy+3y^2}$，求 $u(x,y)$.
}

\oneproblem{
    设位于 $(0,1)$ 的质点 $A$ 对质点 $M$ 的引力大小为 $\dfrac{k}{r^2}$（$k>0$ 为常数，$r$ 为质点 $A$ 与 $M$ 之间的距离），质点 $M$ 沿曲线 $y=\sqrt{2x-x^2}$ 自 $B(2,0)$ 运动到 $O(0,0)$，求在此运动过程中质点 $A$ 对质点 $M$ 的引力所作的功.
}

\oneproblem{
    求 $I=\displaystyle \int_L (e^x\sin y-b(x+y))dx+(e^x\cos y-ax)dy$，其中 $a,b$ 为正的常数，$L$ 为从点 $A(2a,0)$ 沿曲线 $y=\sqrt{2ax-x^2}$ 到点 $O(0,0)$ 的弧.
}

\oneproblem{
    计算曲线积分 $I=\displaystyle \oint_L \dfrac{xdy-ydx}{4x^2+y^2}$，其中 $L$ 是以点 $(1,0)$ 为中心，$R(R>1)$ 为半径的圆周，取逆时针方向.
}

\oneproblem{
    确定常数 $\lambda$，使在右半平面 $x>0$ 上的向量 $\vec{A}(x,y)=2xy(x^4+y^2)^\lambda\vec{i}-x^2(x^4+y^2)^\lambda\vec{j}$ 为某二元函数 $u(x,y)$ 的梯度，并求 $u(x,y)$.
}

\oneproblem{
    设函数 $f(x,y)$ 满足 $\dfrac{\partial f(x,y)}{\partial x}=(2x+1)e^{2x-y}$，且 $f(0,y)=y+1$，$L_t$ 是从点 $(0,0)$ 到点 $(1,t)$ 的光滑曲线，计算曲线积分 $I(t)=\displaystyle \int_{L_t} \dfrac{\partial f(x,y)}{\partial x}dx+\dfrac{\partial f(x,y)}{\partial y}dy$，并求 $I(t)$ 的最小值.
}

\oneproblem{
    设函数 $Q(x,y)$ 在 $xOy$ 面上具有一阶连续偏导数，曲线积分 $\displaystyle \int_L 2xydx+Q(x,y)dy$ 与路径无关，并且对任意 $t$ 恒有
    $
    \int_{(0,0)}^{(t,1)} 2xy\dd x+Q(x,y)dy=\int_{(0,0)}^{(1,t)} 2xy\dd x+Q(x,y)dy,
    $
    求 $Q(x,y)$.
}

\oneproblem{
    设 $D\subset\mathbb{R}^2$ 是有界单连通闭区域，$I(D)=\displaystyle \iint_D (4-x^2-y^2)\dd x \dd y$ 取得最大值的积分区域记作 $D_1$.
    \begin{enumerate}[label=(\arabic*)]
        \item 求 $I(D_1)$ 的值；
        \item 计算 $\displaystyle \int_{\partial D_1} \dfrac{(xe^{x^2+4y^2}+y)dx+(4ye^{x^2+4y^2}-x)dy}{x^2+4y^2}$，其中 $\partial D_1$ 是 $D_1$ 的正向边界.
    \end{enumerate}
}

\oneproblem{
    计算曲线积分 $I=\displaystyle \oint_L \dfrac{udv-vdu}{u^2+v^2}$，其中 $u=ax+by,\ v=cx+dy\ (ad-bc\neq0)$，$L$ 为 $xOy$ 平面上环绕坐标原点的简单闭曲线，方向取逆时针方向.
}

% ======================================================
% 第 77 天
% ======================================================
\setday{77}{积分与路径无关与微分方程求解}

\oneproblem{
    设 $f(x)$ 具有二阶连续导数，$f(0)=0,\ f'(0)=1$，且 $[xy(x+y)-f(x)y]dx+[f'(x)+x^2y]dy=0$ 为一全微分方程，求 $f(x)$ 及此全微分方程的通解.
}

\oneproblem{
    设函数 $f(t)$ 在 $t\neq0$ 时一阶连续可导，且 $f(1)=0$，求函数 $f(x^2-y^2)$，使得曲线积分
    $
    \int_L y[2-f(x^2-y^2)]dx+xf(x^2-y^2)dy
    $
    与路径无关，其中 $L$ 为任一不与直线 $y=\pm x$ 相交的分段光滑曲线.
}

\oneproblem{
    求下列微分方程的通解：
    \begin{enumerate}[label=(\arabic*)]
        \item $(1+xy)ydx+(1-xy)xdy=0$.
        \item $y^2(x-3y)dx+(1-3xy^2)dy=0$.
    \end{enumerate}
}

\oneproblem{
    设曲线积分 $\displaystyle \int_L [f(x)-e^x]\sin ydx-f(x)\cos ydy$ 与路径无关，其中 $f(x)$ 具有一阶连续导数，且 $f(0)=0$，求 $f(x)$ 的表达式.
}

\oneproblem{
    设 $u=u(x)$ 连续可微，$u(2)=1$，且 $\displaystyle \int_L (x+2y)udx+(x+u^3)udy$ 在右半平面上与路径无关，求 $u(x)$.
}

\oneproblem{
    设 $\varphi(x)$ 三次可微，$\varphi(1)=1,\ \varphi'(1)=7$，求 $u(x,y)$，使得 $du=[x^2\varphi'(x)-11x\varphi(x)]dy-32y\varphi(x)dx$.
}

\oneproblem{
    若 $\varphi'(x)$ 连续，$\varphi(\pi)=1$，试确定 $\varphi(x)$，使得积分 $I=\displaystyle \int_A^B [\sin x-\varphi(x)]\dfrac{y}{x}dx+\varphi(x)dy\ (x\neq0)$ 与路线无关，并求 $I=\displaystyle \int_{(1,0)}^{(\pi,\pi)} [\sin x-\varphi(x)]\dfrac{y}{x}dx+\varphi(x)dy$.
}

\oneproblem{
    设 $\varphi(y),\psi(y)$ 二阶连续可导，曲线积分 $\displaystyle \oint_C 2[x\varphi(y)+\psi(y)]dx+[x^2\psi(y)+2xy^2-2x\varphi(y)]dy=0$，这里 $C$ 是平面上任意一条封闭曲线.
    \begin{enumerate}[label=(\arabic*)]
        \item 设 $\varphi(0)=-2,\ \psi(0)=1$，确定 $\varphi(y)$ 与 $\psi(y)$ 的表达式；
        \item 计算从点 $(0,0)$ 到 $\left(\pi,\dfrac{\pi}{2}\right)$ 的任意一条曲线的积分值 $I$.
    \end{enumerate}
}

\oneproblem{
    设 $f(x,y)$ 有连续的偏导数，且 $f(0,0)=0$，满足当 $x^2+y^2\le5$ 时 $|\text{grad }f|\le1$，证明 $|f(1,2)|\le\sqrt{5}$.
}

\oneproblem{
    设函数 $\varphi(y)$ 具有连续导数，在围绕原点的任意分段光滑简单闭曲线 $L$ 上，曲线积分 $\displaystyle \oint_L \dfrac{\varphi(y)dx+2xydy}{2x^2+y^4}$ 的值恒为同一常数.
    \begin{enumerate}[label=(\arabic*)]
        \item 证明：对右半平面 $x>0$ 内的任意分段光滑简单闭曲线 $C$，有 $\displaystyle \oint_C \dfrac{\varphi(y)dx+2xydy}{2x^2+y^4}=0$；
        \item 求函数 $\varphi(y)$ 的表达式；
        \item 设 $C$ 是围绕原点的光滑简单正向闭曲线，求 $\displaystyle \oint_C \dfrac{\varphi(y)dx+2xydy}{2x^2+y^4}$.
    \end{enumerate}
}

\oneproblem{
    设函数 $f(x)$ 在 $(-\infty,+\infty)$ 内具有一阶连续导数，$L$ 是上半平面 $(y>0)$ 内的有向分段光滑曲线，其起点为 $(a,b)$，终点为 $(c,d)$，记
    $
    I=\int_L \dfrac{1}{y}\left[1+y^2f(xy)\right]dx+\dfrac{x}{y^2}\left[y^2f(xy)-1\right]dy.
    $
    \begin{enumerate}[label=(\arabic*)]
        \item 证明曲线积分 $I$ 与路径 $L$ 无关；
        \item 当 $ab=cd$ 时，求 $I$ 的值.
    \end{enumerate}
}

\oneproblem{
    \begin{enumerate}[label=(\arabic*)]
        \item 设 $L$ 为平面内任意一条不经过点 $(x_0,y_0)$ 的正向光滑封闭简单曲线。\\证明：曲线积分 $I=\displaystyle \oint_L \dfrac{(x-x_0)dy-(y-y_0)dx}{c_1(x-x_0)^2+c_2(y-y_0)^2}\ (c_1>0,c_2>0)$ 的值为常数，并求其值.
        \item 求 $I_1=\displaystyle \oint_L \dfrac{xdy-\left(y-\dfrac{1}{2}\right)dx}{x^2+\left(y-\dfrac{1}{2}\right)^2}$，其中 $L$ 从 $A(1,0)$ 经上半单位圆周到 $B(-1,0)$，再经过直线段 $BA$ 回到 $A$ 点.
    \end{enumerate}
}

\oneproblem{
    求极限 $I=\displaystyle \lim_{R\to+\infty} \oint_L \dfrac{xdy-ydx}{(x^2+y^2)^\alpha},\ L:x^2+xy+y^2=R^2$.
}

% ======================================================
% 第 78 天
% ======================================================
\setday{78}{对面积的曲面积分及基本计算方法}

\oneproblem{
    设 $\Sigma: x^2+y^2+z^2=1$，计算对面积的曲面积分 $I=\displaystyle \iint_\Sigma (x^2+2y^2+3z^2)dS$.
}

\oneproblem{
    设曲面 $\Sigma: |x|+|y|+|z|=1$，计算 $I=\displaystyle \oiint_\Sigma (x+|y|)dS $.
}

\oneproblem{
    计算 $I=\displaystyle \oiint_\Sigma (x^2+y^2)dS$，其中 $\Sigma$ 是球面 $x^2+y^2+z^2=2(x+y+z)$.
}

\oneproblem{
    计算曲面积分 $\displaystyle \iint_\Sigma z dS$，其中 $\Sigma$ 为锥面 $z=\sqrt{x^2+y^2}$ 在柱体 $x^2+y^2\le 2x$ 内的部分.
}

\oneproblem{
    设 $S$ 为椭球面 $\dfrac{x^2}{2}+\dfrac{y^2}{2}+z^2=1$ 的上半部分，点 $P(x,y,z)\in S$，$\pi$ 为 $S$ 在点 $P$ 处的切平面，$\rho(x,y,z)$ 为点 $O(0,0,0)$ 到平面 $\pi$ 的距离，求 $\displaystyle \iint_S \dfrac{z dS}{\rho(x,y,z)}$.
}

\oneproblem{
    计算 $I=\displaystyle \iint_\Sigma \dfrac{dS}{\lambda-z}\ (\lambda>R)$，其中 $\Sigma$ 为球面 $\Sigma: x^2+y^2+z^2=R^2$.
}

\oneproblem{
    设 $P$ 为椭球面 $S: x^2+y^2+z^2-yz=1$ 上的动点，若 $S$ 在点 $P$ 处的切平面与 $xOy$ 面垂直，求点 $P$ 的轨迹 $C$，并计算曲面积分 $I=\displaystyle \iint_\Sigma \dfrac{(x+\sqrt{3})|y-2z|}{\sqrt{4+y^2+z^2-4yz}}dS$，其中 $\Sigma$ 是椭球面 $S$ 位于曲线 $C$ 上方的部分.
}

\oneproblem{
    设薄片型物体 $S$ 是圆锥面 $z=\sqrt{x^2+y^2}$ 被柱面 $z^2=2x$ 割下的有限部分，其上任一点密度为 $\mu(x,y,z)=9\sqrt{x^2+y^2+z^2}$。记圆锥与柱面的交线为 $C$：
    \begin{enumerate}[label=(\arabic*)]
        \item 求 $C$ 在 $xOy$ 平面上的投影曲线的方程；
        \item 求 $S$ 的质量 $M$.
    \end{enumerate}
}

\oneproblem{
    证明：$\displaystyle \oiint_\Sigma (x+y+z+\sqrt{3}a)dS\ge12\pi a^3\ (a>0)$，其中 $\Sigma$ 为球面 $(x-a)^2+(y-a)^2+(z-a)^2=a^2$.
}

\oneproblem{
    设函数 $f(x)$ 连续，$a,b,c$ 为常数，$\Sigma$ 是单位球面 $x^2+y^2+z^2=1$。记第一型曲面积分 $I=\displaystyle \iint_\Sigma f(ax+by+cz)dS$。求证：$I=2\pi\displaystyle \int_{-1}^1 f\left(\sqrt{a^2+b^2+c^2}u\right)du$.
}

\oneproblem{
    对于 4 次齐次函数 $f(x,y,z)=a_1x^4+a_2y^4+a_3z^4+3a_4x^2y^2+3a_5y^2z^2+3a_6x^2z^2$，计算曲面积分 $\displaystyle \oiint_\Sigma f(x,y,z)dS$，其中 $\Sigma: x^2+y^2+z^2=1$.
}

\oneproblem{
    计算积分 $I=\displaystyle \int_0^{2\pi} d\phi\int_0^\pi e^{\sin\theta(\cos\phi-\sin\phi)}\sin\theta d\theta$.
}

\oneproblem{
    设曲面 $\Sigma: z^2=x^2+y^2,1\le z\le z$，其面密度为常数 $\rho$。求在原点处的质量为 1 的质点和 $\Sigma$ 之间的引力（记引力常数为 $G$）.
}

\oneproblem{
    已知 $\Sigma$ 是空间曲线 $\begin{cases} x^2+3y^2=1 \\ z=0 \end{cases}$ 绕着 $y$ 旋转而成的椭球面，$S$ 表示曲面 $\Sigma$ 的上半部分（$z\ge0$），$\Pi$ 是椭球面 $S$ 在 $P(x,y,z)$ 点处的切平面，$\rho(x,y,z)$ 是原点到切平面 $\Pi$ 的距离，$\lambda,\mu,\nu$ 表示 $S$ 的外法线的方向余弦。
    \begin{enumerate}[label=(\arabic*)]
        \item 计算 $\displaystyle \iint_S \dfrac{z}{\rho(x,y,z)}dS$；
        \item 计算 $I=\displaystyle \iint_S z(\lambda x+3\mu y+\nu z)dS$，其中 $\Sigma$ 为外侧.
    \end{enumerate}
}

\oneproblem{
    \begin{enumerate}[label=(\arabic*)]
        \item 设 $f(x)$ 在 $[0,1]$ 上可积，证明：$\displaystyle \int_0^{\frac{\pi}{2}} d\varphi\int_0^{\frac{\pi}{2}} f(\cos\varphi\cos\theta)\cos\theta d\theta=\dfrac{\pi}{2}\int_0^1 f(x)dx$.
        \item 计算积分 $I=\displaystyle \int_0^{\frac{\pi}{2}} d\varphi\int_0^{\frac{\pi}{2}} \dfrac{\sin\theta\ln(2-\sin\theta\cos\varphi)}{2-2\sin\theta\cos\varphi+\sin^2\theta\cos^2\varphi}d\theta$.
    \end{enumerate}
}

\oneproblem{
    设函数 $f(t)$ 在 $[0,+\infty)$ 上连续，$\Omega(t)=\{(x,y,z)|x^2+y^2+z^2\le t^2,z\ge0,t>0\}$，$S(t)$ 是 $\Omega(t)$ 的表面，$D(t)$ 是 $\Omega(t)$ 在 $xOy$ 平面的投影区域，$L(t)$ 是 $D(t)$ 的边界曲线。已知对任意 $t\in(0,+\infty)$，恒有
    $
    \oint_{L(t)} f(x^2+y^2)\sqrt{x^2+y^2}ds+\oiint_{S(t)} (x^2+y^2+z^2)dS=\iint_{D(t)} f(x^2+y^2)d\sigma+\iiint_{\Omega(t)} \sqrt{x^2+y^2+z^2}dV
    $
    求 $f(t)$ 的表达式.
}

% ======================================================
% 第 79 天
% ======================================================
\setday{79}{对坐标的曲面积分及基本计算方法}

\oneproblem{
    设 $\Sigma$ 为曲面 $x^2+y^2+4z^2=4\ (z\ge0)$ 的上侧，计算曲面积分 $I=\displaystyle \iint_\Sigma \sqrt{4-x^2-4z^2}\dd x \dd y$.
}

\oneproblem{
    计算 $\displaystyle \iint_\Sigma -ydzdx+(z+1)\dd x \dd y$，其中 $\Sigma$ 是圆柱面 $x^2+y^2=4$ 被平面 $x+z=2$ 和 $z=0$ 所截出部分的外侧.
}

\oneproblem{
    计算曲面积分 $\displaystyle \oiint_S \dfrac{x dydz+z^2 \dd x \dd y}{x^2+y^2+z^2}$，其中 $S$ 是由曲面 $x^2+y^2=R^2$ 及两平面 $z=R,z=-R\ (R>0)$ 所围成立体表面的外侧.
}

\oneproblem{
    设有流速场 $\vec{v}=(0,yz,z^2)$，求穿过曲面 $\Sigma$ 的流量。$\Sigma$ 为柱面 $y^2+z^2=1(z\ge0)$ 被平面 $x=0$ 和 $x=1$ 截下的部分，法向量指向上侧.
}

\oneproblem{
    设 $\Sigma: z=\sqrt{1-x^2-y^2}$，$\gamma$ 是其外法线与 $z$ 轴正向夹成的锐角，计算 $I=\displaystyle \iint_\Sigma z^2\cos\gamma dS$.
}

\oneproblem{
    设 $\Sigma$ 是平面 $P: x+y+z=6$ 围在球面 $(x-1)^2+(y-1)^2+(z-1)^2\le12$ 内部的平面片，方向向下，试用直接法计算对坐标的曲面积分 $I=\displaystyle \iint_\Sigma xdydz+ydzdx+z\dd x \dd y$.
}

\oneproblem{
    设 $\Sigma: \dfrac{x^2}{a^2}+\dfrac{y^2}{b^2}+\dfrac{z^2}{c^2}=1$ 取外侧，计算 $I=\displaystyle \oiint_\Sigma \dfrac{dydz}{x}+\dfrac{dzdx}{y}+\dfrac{\dd x \dd y}{z}$.
}

\oneproblem{
    计算曲面积分 $I=\displaystyle \oiint_\Sigma \dfrac{x dydz+y dzdx+z \dd x \dd y}{(x^2+y^2+z^2)^{\frac{3}{2}}}$，其中 $\Sigma$ 是曲面 $2x^2+2y^2+z^2=4$ 的外侧.
}

\oneproblem{
    设 $S$ 为简单封闭光滑曲面，$P(x,y,z),Q(x,y,z),R(x,y,z)$ 为 $S$ 上的连续函数，\\证明 $\left|\displaystyle \iint_S Pdydz+Qdzdx+R\dd x \dd y\right|\le M\cdot\Delta S$，其中 $M=\max_{(x,y,z)\in S}\sqrt{P^2+Q^2+R^2}$，$\Delta S$ 为 $S$ 的面积.
}

\oneproblem{
    计算 $\displaystyle \iint_\Sigma \dfrac{ax dydz+(z+a)^2 \dd x \dd y}{(x^2+y^2+z^2)^{1/2}}$，其中 $\Sigma$ 为下半球面 $z=-\sqrt{a^2-x^2-y^2}$ 的上侧，$a$ 为大于零的常数.
}

\oneproblem{
    设 $\Sigma$ 是球面 $x^2+y^2+z^2=1$ 的外侧，计算 $I=\displaystyle \iint_\Sigma \dfrac{2dydz}{x\cos^2x}+\dfrac{dzdx}{\cos^2y}-\dfrac{\dd x \dd y}{z\cos^2z}$.
}

\oneproblem{
    设 $\Sigma$ 为曲面 $z=\sqrt{x^2+y^2}\ (1\le x^2+y^2\le4)$ 的下侧，$f(x)$ 为连续函数，计算曲面积分
    $
    I=\iint_\Sigma [xf(xy)+2x-y]\dd y \dd z+[yf(xy)+2y+x]\dd z \dd x+[zf(xy)+z]\dd x \dd y.
    $
}

% ======================================================
% 第 80 天
% ======================================================
\setday{80}{高斯公式 散度}

\oneproblem{
    设 $\Sigma$ 为空间区域 $\{(x,y,z)|x^2+4y^2\le4,0\le z\le2\}$ 表面的外侧，计算曲面积分 $I=\displaystyle \iint_\Sigma x^2dydz+y^2dzdx+z\dd x \dd y$.
}

\oneproblem{
    设数量场 $u=\ln\sqrt{x^2+y^2+z^2}$，计算 $\text{div}(\text{grad }u)$.
}

\oneproblem{
    计算曲面积分 $I=\displaystyle \iint_\Sigma x(8y+1)dydz+2(1-y^2)dzdx-4yz\dd x \dd y$，其中 $\Sigma$ 是由曲线 $\begin{cases} z=\sqrt{y-1} \\ x=0 \end{cases}\ (1\le y\le3)$ 绕 $y$ 旋转一周所成的曲面，它的法向量与 $y$ 轴正向的夹角恒大于 $\dfrac{\pi}{2}$.
}

\oneproblem{
    计算面积分 $\displaystyle \iint_\Sigma (x^3+az^2)dydz+(y^3+ax^2)dzdx+(z^3+ay^2)\dd x \dd y$，其中 $\Sigma$ 为上半球面 $z=\sqrt{a^2-x^2-y^2}$ 的上侧.
}

\oneproblem{
    设 $\Sigma$ 是曲面 $1-\dfrac{z}{5}=\dfrac{(x-2)^2}{16}+\dfrac{(y-1)^2}{9}\ (z\ge0)$ 取上侧，计算 $I=\displaystyle \iint_\Sigma \dfrac{x dydz+y dzdx+z \dd x \dd y}{(x^2+y^2+z^2)^{3/2}}$.
}

\oneproblem{
    设对于半空间 $x>0$ 内任意的光滑有向封闭曲面 $S$，都有 $\displaystyle \oiint_S xf(x)dydz-xyf(x)dzdx-e^{2x}z\dd x \dd y=0$，其中 $f(x)$ 在 $(0,+\infty)$ 内具有连续的一阶导数，且 $\displaystyle \lim_{x\to0^+} f(x)=1$，求 $f(x)$.
}

\oneproblem{
    设 $\Sigma$ 是一个光滑封闭曲面，方向朝外，给定第二型的曲面积分 $I=\displaystyle \iint_\Sigma (x^3-x)dydz+(2y^3-y)dzdx+(3z^3-z)\dd x \dd y$。试确定曲面 $\Sigma$，使得积分 $I$ 的值最小，并求该最小值.
}

\oneproblem{
    若对于 $\mathbb{R}^3$ 中半空间 $\{(x,y,z)\in\mathbb{R}^3|x>0\}$ 内任意有向光滑封闭曲面 $S$，都有
    $
    \iint_S xf'(x)\dd y \dd z+y(xf(x)-f'(x))\dd z \dd x-xz(\sin x+f'(x))\dd x \dd y=0,
    $
    其中 $f$ 在 $(0,+\infty)$ 上二阶导数连续且 $\displaystyle \lim_{x\to0^+} f(x)=\lim_{x\to0^+} f'(x)=0$，求 $f(x)$.
}

\oneproblem{
    设 $\Omega$ 是由光滑的简单封闭曲面 $\Sigma$ 围成的有界闭区域，函数 $f(x,y,z)$ 在 $\Omega$ 上具有连续二阶偏导数，且 $f(x,y,z)|_{(x,y,z)\in\Sigma}=0$。记 $\nabla f$ 为 $f(x,y,z)$ 的梯度，并令 $\Delta f=\dfrac{\partial^2 f}{\partial x^2}+\dfrac{\partial^2 f}{\partial y^2}+\dfrac{\partial^2 f}{\partial z^2}$。证明：对任意常数 $C>0$，恒有
    $
    C\iiint_\Omega f^2\dd x \dd ydz+\dfrac{1}{C}\iiint_\Omega (\Delta f)^2\dd x \dd ydz\ge2\iiint_\Omega |\nabla f|^2\dd x \dd ydz.
    $
}

\oneproblem{
    设函数 $f(x,y,z)$ 在区域 $\Omega=\{(x,y,z)|x^2+y^2+z^2\le1\}$ 上具有连续的二阶偏导数，且满足 $\dfrac{\partial^2 f}{\partial x^2}+\dfrac{\partial^2 f}{\partial y^2}+\dfrac{\partial^2 f}{\partial z^2}=\sqrt{x^2+y^2+z^2}$，计算 $I=\displaystyle \iiint_W \left(x\dfrac{\partial f}{\partial x}+y\dfrac{\partial f}{\partial y}+z\dfrac{\partial f}{\partial z}\right)\dd x \dd ydz$.
}

\oneproblem{
    设 $P(x,y,z)$ 和 $R(x,y,z)$ 在空间上有连续偏导数，设上半球面 $S: z=z_0+\sqrt{r^2-(x-x_0)^2-(y-y_0)^2}$，方向向上，若对任何点 $(x_0,y_0,z_0)$ 和 $r>0$，第二型曲面积分 $\displaystyle \iint_S Pdydz+R\dd x \dd y=0$。证明：$\dfrac{\partial P}{\partial x}\equiv0$.
}

\oneproblem{
    设函数 $f(x)$ 连续可导，$P=Q=R=f((x^2+y^2)z)$，有向曲面 $\Sigma_1$ 是圆柱体 $x^2+y^2\le t^2,0\le z\le1$ 的表面，方向朝外。记第二型曲面积分 $I_t=\displaystyle \iint_{\Sigma_1} Pdydz+Qdzdx+R\dd x \dd y$。求极限 $\displaystyle \lim_{t\to0+}\dfrac{I_t}{t^4}$.
}

\oneproblem{
    \begin{enumerate}[label=(\arabic*)]
        \item 设一球缺高为 $h$，所在球半径为 $R$。证明该球缺的体积为 $\dfrac{\pi}{3}(3R-h)h^2$，球冠的面积为 $2\pi Rh$.
        \item 设球体 $(x-1)^2+(y-1)^2+(z-1)^2\le12$ 被平面 $P: x+y+z=6$ 所截的小球缺为 $\Omega$。记球缺上的球冠为 $\Sigma$，方向指向球外，求第二型曲面积分 $I=\displaystyle \iint_\Sigma xdydz+ydzdx+z\dd x \dd y$.
    \end{enumerate}
}

\oneproblem{
    设 $g(t)=\displaystyle \oiint_{\Sigma_t} zf(x,y)\dd x \dd y$，其中 $\Sigma_t$ 为球面 $x^2+y^2+z^2=t^2\ (t\ge0)$ 的外侧，$f(x,y)$ 在原点处连续。证明：$\displaystyle \lim_{t\to0^+}\dfrac{g(t)}{t^3}=\dfrac{4\pi}{3}f(0,0)$.
}

\oneproblem{
    设 $f(x,y,z)$ 在区域 $\Omega=\{(x,y,z)|x^2+y^2+z^2\le1\}$ 上具有连续的阶偏导数，且满足 $\dfrac{\partial^2 f}{\partial x^2}+\dfrac{\partial^2 f}{\partial y^2}+\dfrac{\partial^2 f}{\partial z^2}=\sqrt{x^2+y^2+z^2}$。计算三重积分 $I=\displaystyle \iiint_\Omega \left(x\dfrac{\partial f}{\partial x}+y\dfrac{\partial f}{\partial y}+z\dfrac{\partial f}{\partial z}\right)dV$.
}